\documentclass[11pt,a4paper]{article}
\usepackage[margin=2.5cm]{geometry}
\usepackage{graphicx}
\usepackage[hidelinks]{hyperref}
\renewcommand{\thesubsection}{\thesection(\alph{subsection})}
\setlength{\parindent}{0pt}
\setlength{\parskip}{5pt}
\title{C1 coursework: strata}
\author{CRSid: replace this \qquad Team: replace this}
\date{Michaelmas 2026}
\begin{document}
\maketitle

% This template is optional. Replace its prompts with your answers.
% Maximum 3000 words; section 4(b) has its own 150-word limit.
% Use the brief for the complete requirements and formulas.
% Keep this file at report/report.tex in your own private submission project.
% Compile from the project root:
% pdflatex -output-directory=report report/report.tex
% Paths in \includegraphics are then relative to the project root.

\section{Implementation of the formulas}
\subsection{JAX functions}
Identify the implementation file and briefly describe relevant choices.
Submit the full code separately.

\subsection{Tests against the reference}
Summarise the three test results against their specified tolerances.
Record the measured errors, angles and test-log path in the JSON.

\section{Cost prediction}
\subsection{Counting constants}
Give one sentence explaining each count. Record all eleven values in the JSON.

\subsection{Predictions at the measured sizes}
Present the predictions in a table, using the JSON values. Label the units.

\subsection{Scaling with the number of individuals}
State the predicted exponents obtained using the supplied formula.

\subsection{Batch size at the ridge point}
State the predicted threshold.

\subsection{Memory capacity}
State the largest number of individuals under the specified capacity rule.

\section{Measurements on the TPU}
Record the device and software version. Use the brief's timing procedure.

\subsection{Formation of the relationship matrix}
Present the timings, fitted exponent and buffer sizes. Compare the fitted
exponent with Question 2(c), and report the host-CPU time and speed-up.

\subsection{Eigenvector calculations}
Present the timings, their ratio and the principal angles for the specified
two-column comparison.

\subsection{Batched association scan}
Include the plot of execution time against batch size, with a caption and
labelled axes. Compare the first tested crossing with the predicted threshold.
% \begin{figure}[htbp]
% \centering
% \includegraphics[width=0.8\linewidth]{figures/scan.pdf}
% \caption{Replace with a caption describing the plotted results.}
% \end{figure}

\section{Roofline analysis}
\subsection{Measured rates on the roofline}
Include the labelled roofline plot and a caption.
% \begin{figure}[htbp]
% \centering
% \includegraphics[width=0.8\linewidth]{figures/roofline.pdf}
% \caption{Replace with a caption describing the plotted results.}
% \end{figure}

\subsection{Interpretation}
In at most 150 words, explain which kernels approach the roofline and which
fall below it. Relate the differences to their calculations and memory traffic.

\end{document}
